% ------------------------------------------------------------
\begin{frame}[fragile]
  \frametitle{Docker -- Running MySQL Without XAMPP}

  \begin{block}{What Is Docker?}
    \textbf{Docker} packages an application and all its dependencies into a
    self-contained unit called a \textbf{container}. A container runs identically
    on any machine -- no ``it works on my laptop'' problems.\\[0.3em]
    For databases, Docker means: spin up a MySQL server in \emph{one command},
    destroy it just as easily, and run multiple versions side by side.
  \end{block}

  \begin{block}{Starting a MySQL container}
\begin{lstlisting}[style=sqlstyle, language=bash]
# Pull the official MySQL image and start a container
docker run --name my-mysql \
  -e MYSQL_ROOT_PASSWORD=secret \
  -e MYSQL_DATABASE=supermarket \
  -p 3306:3306 \
  -d mysql:8.0

# Connect with the MySQL CLI inside the container
docker exec -it my-mysql mysql -u root -psecret

# Stop and remove the container when done
docker stop my-mysql && docker rm my-mysql
\end{lstlisting}
  \end{block}
\end{frame}

% ------------------------------------------------------------
\begin{frame}[fragile]
  \frametitle{Docker Compose -- Multi-Container Setup}

  Real projects often need both a database \emph{and} an application server.
  \textbf{Docker Compose} describes the entire stack in a single YAML file.

  \begin{block}{docker-compose.yml example}
\begin{lstlisting}[style=sqlstyle]
version: "3.9"
services:
  db:
    image: mysql:8.0
    environment:
      MYSQL_ROOT_PASSWORD: secret
      MYSQL_DATABASE: supermarket
    ports:
      - "3306:3306"
    volumes:
      - db_data:/var/lib/mysql   # persist data across restarts

  adminer:
    image: adminer              # lightweight phpMyAdmin alternative
    ports:
      - "8080:8080"

volumes:
  db_data:
\end{lstlisting}
  \end{block}

  \begin{examples}{Usage}
    \texttt{docker compose up -d} starts all services.
    \texttt{docker compose down} stops them.
    Connect DBeaver to \texttt{localhost:3306} as usual.
  \end{examples}
\end{frame}

% ------------------------------------------------------------
\begin{frame}
  \frametitle{Tool Comparison}

  \begin{center}
  \renewcommand{\arraystretch}{1.3}
  \small
  \begin{tabular}{@{}lllll@{}}
    \toprule
    \textbf{Tool}      & \textbf{Type}      & \textbf{DB Support} & \textbf{Best For} \\
    \midrule
    XAMPP              & Local server bundle & MySQL/MariaDB       & Quick local setup \\
    phpMyAdmin         & Browser GUI         & MySQL/MariaDB       & Simple admin tasks \\
    MySQL Workbench    & Desktop GUI         & MySQL only          & ERM design, DDL \\
    DBeaver            & Desktop GUI         & Any (via JDBC)      & Daily SQL work \\
    Docker             & Container runtime   & Any image           & Reproducible envs \\
    \bottomrule
  \end{tabular}
  \end{center}

  \begin{block}{Recommendation for this course}
    \begin{itemize}
      \item \textbf{Database server}: XAMPP (Windows/macOS) or Docker (any OS)
      \item \textbf{Schema design}: MySQL Workbench (ERM forward-engineering)
      \item \textbf{Daily SQL queries}: DBeaver
    \end{itemize}
  \end{block}
\end{frame}
